\section{Scope, main results, and mathematical conventions}
\label{scope:section}

The objective is to obtain exact identities that survive the passage
between spectral parameters, endpoint singularities, and ordinary
special functions. Four aspects of the existing research collection
are particularly productive: centered harmonic expansions,
cotangent--Hurwitz transforms, entire relative harmonic interpolants,
and holomorphic remainders of multivariate zeta functions. Each already
has a substantial foundation in the manuscript or the incoming reports.
The present article extends that foundation at specific open points.

\subsection{The inspected baseline and the status of the questions}

The source baseline is the repository revision \snapshot,
dated 11 October 2026 at 00:52:39 UTC. We inspected the relevant
canonical manuscript chapters and the nine incoming archives listed
in the accompanying source manifest. The most direct starting points
are \emph{Exact Resonant Identities} \cite{ghp:ExactResonant},
\emph{Independent Orders and Exact Reductions} \cite{qtr:Independent},
and \emph{Cubic Tornheim Derivatives and Harmonic Stieltjes
Constants} \cite{source:Cubic}. The original manuscript remains
the reference for its established normalization and conjecture labels
\cite{source:manuscript}.

This is a targeted continuation and audit of those areas, not a claim
to have checked every line of all manuscript volumes. In particular,
the short $S_6$ reduction and the revised $S_8$ reduction remain
conjectures. Previously rejected candidate vectors and previously
reported corrections are distinguished from current candidates in
Section~\ref{audit:section}.

The following table states the exact scope of the main advances.
Resolution refers to the indicated mathematical question, with the
domains and qualifications in the cited theorem.

\begin{table}[htbp]
\centering
\small
\begin{tabularx}{\textwidth}{@{}p{0.23\textwidth}Y p{0.19\textwidth}@{}}
\toprule
Source question or direction & Result proved here & Location\\
\midrule
Exact Resonant Q11: generalized harmonic cancellation
& Necessary and sufficient centered involution criterion for all
nonpositive even spectral points.
& Theorem~\ref{ghp:thm:classification}\\
\addlinespace
Explicit centered values beyond raw $H_n$ powers
& Every trigamma-square even moment, its ordinary convergent
subtracted sum, and a convergent polylogarithmic generator.
& Theorems~\ref{ghp:thm:square}, \ref{ghp:thm:generator}\\
\addlinespace
Exact Resonant Q8: a nonpolynomial cotangent family
& The complete rational sector bounded at infinity, with the
all-index endpoint correction.
& Theorem~\ref{res:main}; Corollary~\ref{res:repeated}\\
\addlinespace
Exact Resonant Q9: normalized primitive families
& The resolvent transform of every balanced Gamma primitive and
the continuous-order conversion formula.
& Theorem~\ref{res:continuous}; Corollary~\ref{res:balancedcor}\\
\addlinespace
Independent Orders: nonpositive-index reductions
& Arbitrary positions and depths, a canonical finite normal form,
and an analytic zero-decoration generator.
& Theorems~\ref{harm:local}--\ref{harm:Lerch}\\
\addlinespace
Independent Orders and Cubic Harmonic: quartic rays
& All quartic directions in three coordinates; cyclic directions
in one diagonal coordinate.
& Theorem~\ref{qtr:quartic}; equation~\eqref{qtr:cyclicquartic}\\
\addlinespace
Higher directional relations
& Exact formal kernel and cyclic-image dimensions for the named
relations, and a fifth-order obstruction.
& Theorem~\ref{qtr:kernel}; Corollary~\ref{qtr:cyclicimage}\\
\bottomrule
\end{tabularx}
\caption{Results relative to the pinned source questions. The primitive
and higher-order statements resolve the specified families, rather than
every possible extension proposed in the source.}
\label{scope:table}
\end{table}

\subsection{A selection of exact identities}

Three compact examples illustrate the different kinds of conclusion.
The centered square identity
\begin{equation}\label{scope:sample-square}
 \sum_{n=0}^{\infty}
 \left((n+\tfrac12)^2\psi'(n+1)^2-1\right)
       =-\frac14\zeta(3)-\frac16
\end{equation}
is an ordinary absolutely convergent series. It is the first
nontrivial member of a complete all-even family, rather than a
numerically recognized isolated value.

The pointwise endpoint finite part
\begin{equation}\label{scope:sample-resolvent}
 \operatorname{FP}_x\int_0^1
 \frac{\psi'(x)}{\pi\cot(\pi x)-i\pi}\,dx
       =1-\gamma-\log(2\pi)+\frac{i\pi}{2}
\end{equation}
has a boundary contribution equal to $1$. Analytically continuing an
integrated spectral family and then substituting the integer order
without that contribution gives a different value. The all-index
theorem locates this difference before any expansion.

For the Tornheim restriction $W_{a,b,c}(t)=T(at,bt,ct)$ and
$\Lambda=W_{1,1,1}^{(4)}(0)$, a cyclic combination gives
\begin{align}\label{scope:sample-quartic}
 &W_{1,2,3}^{(4)}(0)+W_{2,3,1}^{(4)}(0)
       +W_{3,1,2}^{(4)}(0)-36\Lambda\notag\\
 &\hspace{13mm}=-184\log(2\pi)\zeta'''(0)
                +156\zeta''(0)^2-386\zeta''''(0).
\end{align}
The cancellation is an exact polynomial identity between directional
coefficients. It does not require an assumption about relations among
the individual constants.

\subsection{Conventions and the meaning of a reduction}

We use
\[
 \zeta(1+t,x)=\frac1t+
      \sum_{m\ge0}\frac{(-1)^m}{m!}\gamma_m(x)t^m,\qquad
 \gamma_m=\gamma_m(1),\qquad \gamma=\gamma_0,
\]
and $\psi=\Gamma'/\Gamma$. Primes on a one-variable zeta function, or
$\zeta^{(j)}(s,x)$ with the first argument displayed, mean derivatives
in the spectral argument. Superscripts on $\psi$ and $\gamma_m(x)$
mean derivatives in $x$. When confusion is possible, a square
superscript $[j]$ denotes a polylogarithmic order derivative:
\[
 \Li_s^{[j]}(z)=\partial_s^j\Li_s(z).
\]
All logarithms of positive real quantities are real. Complex powers
use the specified branch; Hurwitz endpoints are initially in the
right half-plane. Bernoulli numbers satisfy $B_1=-1/2$.
All branches required by the resolvent formulas are given explicitly
in Section~\ref{res:section}.

The term \emph{reduction} is always relative to a stated target class.
For the harmonic deletion theorem it means a finite linear combination
with polynomial endpoint coefficients and only the surviving free
orders. For a rational cotangent kernel it means finitely many
polylogarithmic order derivatives at explicit arguments. For the
quartic Tornheim layer it means three explicitly normalized
coordinates, two of which are supplied with ordinary convergent
representatives. It does not silently mean reduction to a conjectured
algebra of familiar constants.

There are three distinct analytic operations in the proofs.
A meromorphic continuation is uniquely determined from a convergence
domain. A coordinate finite part extracts the constant term of a
cutoff integral in its stated endpoint variable. A directional
restriction of a meromorphic function is formed before taking its
one-variable limit. These operations sometimes agree and sometimes
differ by explicit terms. The article proves the needed compatibility
at each use.

\subsection{Proof inputs and attribution}

The sufficient half-shift parity argument is elementary asymptotic
algebra. Its converse uses the formal-Laurent version of the
Adamczewski--Dreyfus--Hardouin theorem
\cite[Theorem~1.2]{ghp:ADH}. This use is spelled out, including a
proof that the relevant formal trigamma series is not rational.
It supplies differential algebraic independence of functions and
formal series; it says nothing about algebraic independence of
Euler's constant or zeta values.

The Mellin, Fourier, Bernoulli, Gamma, and polylogarithm expansions
used below are classical \cite{DLMF}. The relative harmonic
interpolant and the Tornheim polar normalization are inherited from
the pinned incoming reports; the underlying one-endpoint
interpolant is due to Ihara, Nakamura, and Yamamoto \cite{harm:INY}.
The continuous polygamma function and balanced negapolygamma functions
are the established Espinosa--Moll functions \cite{EM}.
Polynomial Hurwitz and Stieltjes primitives are also established
in Espinosa--Moll \cite{harm:EM} and the canonical manuscript.
Our fully indexed arbitrary-endpoint version is included to make the
new finite reduction calculus directly usable.

The centered trigamma-square family should also be distinguished
from the weighted squares of \emph{differences} of trigamma functions
already evaluated in the incoming mixed-spectral report
\cite{source:Mixed}. Both are exact identities; their summands,
parameters, and generating mechanisms are different.

All results described as developments are developments relative to
the inspected source baseline. The article makes no literature-wide
priority claim. Proofs are analytic and algebraic; the accompanying
programs provide finite exact checks and numerical diagnostics, not
proof-assistant verification of the analytic theorems.
